\Article{3}
Sind die Größen \(v,v',v''\) etc. völlig voneinander unabhängig, so kann die Unbekannte nur auf eine einzige Weise durch sie bestimmt werden; daher lässt sich jene Unsicherheit dann auf keine Weise vermeiden oder verringern, und hinsichtlich des aus den Beobachtungen abzuleitenden Werts der Unbekannten bleibt nichts der freien Wahl überlassen.

Ganz anders verhält es sich jedoch, wenn zwischen den Größen \(v,v',v''\) etc. eine gegenseitige Abhängigkeit besteht, die, wie wir voraussetzen wollen, durch \(\sigma\) Bedingungsgleichungen
\[
X=0,\quad Y=0,\quad Z=0\ \text{etc.}
\]
ausgedrückt wird, wobei \(X,Y,Z\) etc. gegebene Funktionen der Unbestimmten \(v,v',v''\) etc. bezeichnen. In diesem Fall lässt sich unsere Unbekannte auf unendlich viele verschiedene Arten durch Kombinationen der Größen \(v,v',v''\) etc. bestimmen, da offenbar anstelle der Funktion \(u\) jede andere Funktion \(U\) genommen werden kann, die so beschaffen ist, dass \(U-u\) identisch verschwindet, wenn \(X=0,\ Y=0,\ Z=0\) etc. gesetzt wird.

Bei der Anwendung auf einen bestimmten Fall ergäbe sich hieraus hinsichtlich des Werts der Unbekannten kein Unterschied, wenn die Beobachtungen absolute Genauigkeit hätten; soweit sie aber Fehlern unterworfen bleiben, wird offenbar im Allgemeinen eine andere Kombination einen anderen Wert der Unbekannten liefern. Beispielsweise haben wir nun anstelle des Fehlers
\[
le+l'e'+l''e''+\text{etc.}
\]
den die Funktion \(u\) hervorgebracht hatte,
\[
Le+L'e'+L''e''+\text{etc.}
\]
wenn wir die Funktion \(U\) nehmen und die Werte der Ableitungen \(\frac{\mathrm dU}{\mathrm dv},\frac{\mathrm dU}{\mathrm dv'},\frac{\mathrm dU}{\mathrm dv''}\) etc. jeweils mit \(L,L',L''\) etc. bezeichnen. Obwohl wir die Fehler selbst nicht angeben können, lassen sich dennoch die bei den verschiedenen Kombinationen von Beobachtungen zu be%
\SourceLeaf{065}%
fürchtenden mittleren Fehler miteinander vergleichen; die beste Kombination ist diejenige, bei der dieser mittlere Fehler möglichst klein wird. Da er gleich dem folgenden Ausdruck wird:
\[
=\BareRoot{(LLmm+L'L'm'm'+L''L''m''m''+\text{etc.})}
\]
muss darauf hingearbeitet werden, dass die Summe \(LLmm+L'L'm'm'+L''L''m''m''+\text{etc.}\) ihren kleinsten Wert annimmt.

\Article{4}
Da die unendliche Vielfalt der Funktionen \(U\), die nach der im vorhergehenden Artikel ausgesprochenen Bedingung an die Stelle von \(u\) treten können, hier nur insofern zu betrachten ist, als sich daraus verschiedene Systeme von Werten der Koeffizienten \(L,L',L''\) etc. ergeben, ist vor allem der Zusammenhang zu untersuchen, der zwischen sämtlichen zulässigen Systemen bestehen muss. Bezeichnen wir die festen Werte der partiellen Ableitungen
\begingroup\renewcommand{\arraystretch}{2}
\[
\begin{array}{ccc l}
\dfrac{\mathrm dX}{\mathrm dv},&\dfrac{\mathrm dX}{\mathrm dv'},&\dfrac{\mathrm dX}{\mathrm dv''}&\text{etc.}\\[5pt]
\dfrac{\mathrm dY}{\mathrm dv},&\dfrac{\mathrm dY}{\mathrm dv'},&\dfrac{\mathrm dY}{\mathrm dv''}&\text{etc.}\\[5pt]
\dfrac{\mathrm dZ}{\mathrm dv},&\dfrac{\mathrm dZ}{\mathrm dv'},&\dfrac{\mathrm dZ}{\mathrm dv''}&\text{etc. etc.}
\end{array}
\]
\endgroup
die sie annehmen, wenn man für \(v,v',v''\) etc. die wahren Werte einsetzt, jeweils mit
\[
\begin{array}{ccc l}
a,&a',&a''&\text{etc.}\\
b,&b',&b''&\text{etc.}\\
c,&c',&c''&\text{etc. etc.}
\end{array}
\]
so ist klar: Denkt man sich den Größen \(v,v',v''\) etc. solche Zuwächse \(\mathrm dv,\mathrm dv',\mathrm dv''\) etc. erteilt, durch die \(X,Y,Z\) etc. unverändert bleiben, sodass jede \(=0\) bleibt, das heißt Zuwächse, die den Gleichungen
\[
\begin{aligned}
0&=a\mathrm dv+a'\mathrm dv'+a''\mathrm dv''+\text{etc.}\\
0&=b\mathrm dv+b'\mathrm dv'+b''\mathrm dv''+\text{etc.}\\
0&=c\mathrm dv+c'\mathrm dv'+c''\mathrm dv''+\text{etc. etc.}
\end{aligned}
\]
genügen, so muss auch \(u-U\) unverändert bleiben, und folglich gilt
\[
0=(l-L)\mathrm dv+(l'-L')\mathrm dv'+(l''-L'')\mathrm dv''+\text{etc.}
\]
\SourceLeaf{066}%
Hieraus wird leicht geschlossen, dass die Koeffizienten \(L,L',L''\) etc. durch Formeln der folgenden Art dargestellt werden müssen:
\[
\begin{aligned}
L&=l+ax+by+cz+\text{etc.}\\
L'&=l'+a'x+b'y+c'z+\text{etc.}\\
L''&=l''+a''x+b''y+c''z+\text{etc.}
\end{aligned}
\]
etc., wobei \(x,y,z\) etc. feste Multiplikatoren bezeichnen. Umgekehrt ist klar, dass bei beliebiger Annahme eines Systems fester Multiplikatoren \(x,y,z\) etc. stets eine solche Funktion \(U\) angegeben werden kann, der Werte von \(L,L',L''\) etc. entsprechen, die mit diesen Gleichungen übereinstimmen, und die nach der im vorhergehenden Artikel ausgesprochenen Bedingung an die Stelle von \(u\) treten kann; dies lässt sich sogar auf unendlich viele verschiedene Arten erreichen. Am einfachsten ist es, \(U=u+xX+yY+zZ+\text{etc.}\) zu setzen; allgemeiner darf man \(U=u+xX+yY+zZ+\text{etc.}+u'\) setzen, wobei \(u'\) eine solche Funktion der Unbestimmten \(v,v',v''\) etc. bezeichnet, die für \(X=0,\ Y=0,\ Z=0\) etc. stets verschwindet und deren Wert in dem vorliegenden bestimmten Fall ein Maximum oder Minimum wird. Für unseren Zweck ergibt sich daraus jedoch kein Unterschied.

\clearpage
